\documentclass[11pt]{article}
\usepackage{amsmath, amssymb, geometry, graphicx, booktabs, enumerate}
\geometry{letterpaper, margin=1in}

\title{\textbf{Homework 16.1: Vector Fields\\Comprehensive Problem Sets}}
\author{Math Study Guide}
\date{}

\begin{document}

\maketitle

\section*{Concept Checklist}
Based on the provided source document, the following core concepts, traps, and techniques are covered in this problem set:
\begin{enumerate}
    \item \textbf{Direct Evaluation:} Plugging coordinates into 2D/3D vector fields.
    \item \textbf{Vector Field Sketching \& Matching (2D/3D):} Utilizing the "Axis Trick", "Component Zero Trick", and analyzing magnitude/direction to visualize radial, rotational, or constant flows.
    \item \textbf{Gradient Vector Fields ($\nabla f$):} Calculating partial derivatives, applying the Chain Rule and the Product Rule strictly.
    \item \textbf{Linear Approximation / Kinematics (Euler's Method):} Treating the vector field as velocity to approximate new positions ($\Delta \mathbf{r} \approx \mathbf{V} \Delta t$).
    \item \textbf{Dot Product for Orthogonality:} Using the dot product geometrically to prove fields are purely rotational ($\mathbf{u} \cdot \mathbf{v} = 0$).
    \item \textbf{Conservative Vector Fields \& Cross-Partial Test:} Verifying if a field is the gradient of a potential function by checking $\frac{\partial P}{\partial y} = \frac{\partial Q}{\partial x}$.
    \item \textbf{Diagnostic Error Analysis ("Find the Flaw"):} Identifying common algebraic and conceptual traps like incorrect partial derivative rules or incorrectly applying Euler's method formulas.
\end{enumerate}

\hrule\vspace{0.5cm}

\section*{Homework Problems and Solutions}

\subsection*{Part 1: Direct Evaluation}

\textbf{Problem 1.} Evaluate the vector field $\mathbf{F}(x, y) = \langle x^2 - y^2, 2xy \rangle$ at the point $(-2, 3)$. \\
\textbf{Solution:} 
Substitute $x = -2$ and $y = 3$ into the components:
\begin{align*}
    P(-2, 3) &= (-2)^2 - (3)^2 = 4 - 9 = -5 \\
    Q(-2, 3) &= 2(-2)(3) = -12
\end{align*}
\textbf{Final Answer:} $\mathbf{F}(-2, 3) = \langle -5, -12 \rangle$.

\vspace{0.3cm}
\textbf{Problem 2.} Evaluate the 3D vector field $\mathbf{F}(x, y, z) = x\mathbf{i} - z\mathbf{j} + y\mathbf{k}$ at the point $(0, -4, 5)$. \\
\textbf{Solution:}
The field in component form is $\langle x, -z, y \rangle$. Substituting $x = 0, y = -4, z = 5$:
\begin{align*}
    \mathbf{F}(0, -4, 5) &= 0\mathbf{i} - (5)\mathbf{j} + (-4)\mathbf{k} = \langle 0, -5, -4 \rangle
\end{align*}
\textbf{Final Answer:} $\langle 0, -5, -4 \rangle$.

\vspace{0.3cm}
\textbf{Problem 3.} Find the magnitude of the vector field $\mathbf{F}(x, y) = \langle 4x, 3y \rangle$ exactly at the coordinate $(2, -1)$. \\
\textbf{Solution:}
First, evaluate the vector at the point: $\mathbf{F}(2, -1) = \langle 4(2), 3(-1) \rangle = \langle 8, -3 \rangle$. 
Next, compute the magnitude:
$$|\mathbf{F}| = \sqrt{8^2 + (-3)^2} = \sqrt{64 + 9} = \sqrt{73}$$
\textbf{Final Answer:} $\sqrt{73}$.

\subsection*{Part 2: Vector Field Sketching \& Matching (2D and 3D)}

\textbf{Problem 4.} Describe the visual behavior of the vector field $\mathbf{F}(x, y) = \langle -x, -y \rangle$. Does it point inward or outward, and what happens to the arrow lengths as you move away from the origin? \\
\textbf{Solution:}
Using test points: At $(1,0)$, $\mathbf{F} = \langle -1, 0 \rangle$. At $(0,1)$, $\mathbf{F} = \langle 0, -1 \rangle$. The vectors always point in the exact opposite direction of the position vector $\mathbf{r} = \langle x, y \rangle$. 
The magnitude is $|\mathbf{F}| = \sqrt{(-x)^2 + (-y)^2} = \sqrt{x^2+y^2}$. 
\textbf{Final Answer:} The arrows point radially \textit{inward} toward the origin. Their lengths \textit{increase} as you move further away from the origin.

\vspace{0.3cm}
\textbf{Problem 5.} Apply the "Axis Trick" to determine the behavior of $\mathbf{F}(x,y) = \langle y^2, x^2 \rangle$ on the x-axis and y-axis. \\
\textbf{Solution:}
On the positive x-axis (where $y=0, x>0$): $\mathbf{F}(x, 0) = \langle 0, x^2 \rangle$. The arrows point purely straight up (positive j direction).
On the positive y-axis (where $x=0, y>0$): $\mathbf{F}(0, y) = \langle y^2, 0 \rangle$. The arrows point purely to the right (positive i direction).
\textbf{Final Answer:} On the x-axis, arrows point strictly vertical (up). On the y-axis, arrows point strictly horizontal (right).

\vspace{0.3cm}
\textbf{Problem 6.} A vector field is given by $\mathbf{F}(x, y) = \frac{\langle x, y \rangle}{x^2 + y^2}$. Show algebraically what happens to the length of the arrows as you move further from the origin. \\
\textbf{Solution:}
Calculate the magnitude $|\mathbf{F}|$:
$$|\mathbf{F}| = \sqrt{ \left(\frac{x}{x^2+y^2}\right)^2 + \left(\frac{y}{x^2+y^2}\right)^2 } = \sqrt{\frac{x^2 + y^2}{(x^2+y^2)^2}} = \sqrt{\frac{1}{x^2+y^2}} = \frac{1}{\sqrt{x^2+y^2}}$$
\textbf{Final Answer:} As the distance $\sqrt{x^2+y^2}$ increases, the magnitude $\frac{1}{\sqrt{x^2+y^2}}$ decreases. The arrows shrink as they get further from the origin.

\vspace{0.3cm}
\textbf{Problem 7.} Describe the flow of $\mathbf{F}(x,y) = \langle 0, x \rangle$. What does the "Component Zero Trick" tell you about the arrows? \\
\textbf{Solution:}
Since the x-component $P = 0$ everywhere, every single arrow in this field is perfectly vertical. For $x > 0$ (right half plane), arrows point up. For $x < 0$ (left half plane), arrows point down.
\textbf{Final Answer:} The field consists entirely of vertical arrows. This is a shear flow pointing up on the right side of the y-axis and down on the left side.

\vspace{0.3cm}
\textbf{Problem 8.} Which vector field represents a constant counter-clockwise rotation where the arrow lengths are always 1? 
A) $\langle -y, x \rangle$ \quad B) $\frac{\langle -y, x \rangle}{\sqrt{x^2+y^2}}$ \quad C) $\langle y, -x \rangle$ \\
\textbf{Solution:}
Counter-clockwise rotational flow dictates that at $(1,0)$, the arrow points up $\langle 0,1 \rangle$. Both A and B do this. However, we need a magnitude of exactly 1 everywhere.
Magnitude of A: $\sqrt{(-y)^2+x^2} = \sqrt{x^2+y^2}$ (varies).
Magnitude of B: $\sqrt{ \frac{y^2}{x^2+y^2} + \frac{x^2}{x^2+y^2} } = \sqrt{1} = 1$.
\textbf{Final Answer:} B.

\vspace{0.3cm}
\textbf{Problem 9.} Describe the 3D vector field $\mathbf{F}(x, y, z) = \langle 0, 0, z^2 \rangle$. \\
\textbf{Solution:}
The $x$ and $y$ components are zero, meaning arrows never point left/right or forward/back. The $z$ component is $z^2$, which is always positive (except at $z=0$).
\textbf{Final Answer:} All arrows point perfectly straight up (positive z-direction) everywhere in 3D space, and they get longer as you move further above or below the $z=0$ plane.

\vspace{0.3cm}
\textbf{Problem 10.} How does the 3D vector field $\mathbf{F}(x,y,z) = \langle x, y, -2 \rangle$ behave? \\
\textbf{Solution:}
In the xy-plane, $\langle x, y \rangle$ represents a radially outward explosion from the z-axis. The z-component is a constant $-2$. 
\textbf{Final Answer:} The arrows point radially outwards away from the z-axis, while simultaneously all pointing diagonally downwards at a constant vertical rate.

\vspace{0.3cm}
\textbf{Problem 11.} Use test points to determine the direction of $\mathbf{F}(x,y,z) = \langle -y, x, z \rangle$ at $(1, 0, 1)$ and $(0, 1, 1)$. \\
\textbf{Solution:}
At $(1, 0, 1) \implies \mathbf{F} = \langle 0, 1, 1 \rangle$. (Points in the +y and +z direction).
At $(0, 1, 1) \implies \mathbf{F} = \langle -1, 0, 1 \rangle$. (Points in the -x and +z direction).
\textbf{Final Answer:} This creates an upward spiraling (tornado-like) field.

\subsection*{Part 3: Gradient Vector Fields}

\textbf{Problem 12.} Find the gradient vector field $\nabla f$ of $f(x, y) = x^4 y^3 - 2x + 5y$. \\
\textbf{Solution:}
Find the partial derivatives $f_x$ and $f_y$.
$f_x = \frac{\partial}{\partial x}(x^4 y^3 - 2x + 5y) = 4x^3 y^3 - 2$.
$f_y = \frac{\partial}{\partial y}(x^4 y^3 - 2x + 5y) = 3x^4 y^2 + 5$.
\textbf{Final Answer:} $\nabla f(x, y) = \langle 4x^3 y^3 - 2, 3x^4 y^2 + 5 \rangle$.

\vspace{0.3cm}
\textbf{Problem 13.} Find the gradient vector field $\nabla f$ of $f(x, y) = 3x \sin(y^2)$. \\
\textbf{Solution:}
Treating $y$ as constant: $f_x = 3 \sin(y^2)$.
Treating $x$ as constant, use the chain rule for $\sin(y^2)$: $f_y = 3x \cdot \cos(y^2) \cdot (2y) = 6xy \cos(y^2)$.
\textbf{Final Answer:} $\nabla f(x, y) = \langle 3 \sin(y^2), 6xy \cos(y^2) \rangle$.

\vspace{0.3cm}
\textbf{Problem 14.} Find the gradient vector field $\nabla f$ of $f(x, y) = e^{x^2 y}$. \\
\textbf{Solution:}
Apply the chain rule for $e^u$, where $u = x^2 y$.
$f_x = e^{x^2 y} \cdot \frac{\partial}{\partial x}(x^2 y) = 2xy e^{x^2 y}$.
$f_y = e^{x^2 y} \cdot \frac{\partial}{\partial y}(x^2 y) = x^2 e^{x^2 y}$.
\textbf{Final Answer:} $\nabla f(x, y) = \langle 2xy e^{x^2 y}, x^2 e^{x^2 y} \rangle$.

\vspace{0.3cm}
\textbf{Problem 15.} (Product Rule Trap) Find the gradient vector field $\nabla f$ of $f(x, y) = x^3 y^2 e^{xy}$. \\
\textbf{Solution:}
For $f_x$, both $x^3$ and $e^{xy}$ contain $x$, requiring the Product Rule. Treat $y$ as a constant.
$f_x = (3x^2 y^2)(e^{xy}) + (x^3 y^2)(y e^{xy}) = 3x^2 y^2 e^{xy} + x^3 y^3 e^{xy}$.
For $f_y$, use the Product Rule similarly.
$f_y = (2x^3 y)(e^{xy}) + (x^3 y^2)(x e^{xy}) = 2x^3 y e^{xy} + x^4 y^2 e^{xy}$.
\textbf{Final Answer:} $\nabla f(x, y) = \langle 3x^2 y^2 e^{xy} + x^3 y^3 e^{xy}, 2x^3 y e^{xy} + x^4 y^2 e^{xy} \rangle$.

\vspace{0.3cm}
\textbf{Problem 16.} Find $\nabla f$ for the 3D scalar function $f(x, y, z) = x^2 \cos(yz)$. \\
\textbf{Solution:}
Compute $f_x$, $f_y$, and $f_z$.
$f_x = 2x \cos(yz)$.
$f_y = x^2 \cdot (-\sin(yz)) \cdot z = -x^2 z \sin(yz)$.
$f_z = x^2 \cdot (-\sin(yz)) \cdot y = -x^2 y \sin(yz)$.
\textbf{Final Answer:} $\nabla f(x, y, z) = \langle 2x \cos(yz), -x^2 z \sin(yz), -x^2 y \sin(yz) \rangle$.

\vspace{0.3cm}
\textbf{Problem 17.} Find $\nabla f$ of $f(x, y) = \ln(x^2 + 3y^4)$. \\
\textbf{Solution:}
Use the chain rule for natural log: $\frac{\partial}{\partial x} \ln(u) = \frac{u'}{u}$.
$f_x = \frac{1}{x^2 + 3y^4} \cdot (2x) = \frac{2x}{x^2 + 3y^4}$.
$f_y = \frac{1}{x^2 + 3y^4} \cdot (12y^3) = \frac{12y^3}{x^2 + 3y^4}$.
\textbf{Final Answer:} $\nabla f(x, y) = \left\langle \frac{2x}{x^2 + 3y^4}, \frac{12y^3}{x^2 + 3y^4} \right\rangle$.

\vspace{0.3cm}
\textbf{Problem 18.} Find the gradient vector field of $f(s, t) = \sqrt{5s^3 + 2t}$. \\
\textbf{Solution:}
Rewrite as $f(s,t) = (5s^3 + 2t)^{1/2}$. Use the generalized power rule.
$f_s = \frac{1}{2}(5s^3 + 2t)^{-1/2} \cdot (15s^2) = \frac{15s^2}{2\sqrt{5s^3 + 2t}}$.
$f_t = \frac{1}{2}(5s^3 + 2t)^{-1/2} \cdot (2) = \frac{1}{\sqrt{5s^3 + 2t}}$.
\textbf{Final Answer:} $\nabla f(s, t) = \left\langle \frac{15s^2}{2\sqrt{5s^3 + 2t}}, \frac{1}{\sqrt{5s^3 + 2t}} \right\rangle$.

\subsection*{Part 4: Linear Approximation / Kinematics}

\textbf{Problem 19.} A particle moves in a velocity field $\mathbf{V}(x, y) = \langle y^2, -x \rangle$. If it is at position $(2, 3)$ at time $t = 0$, estimate its location at time $t = 0.1$. \\
\textbf{Solution:}
Use Euler's Method: $\text{New Position} \approx \text{Old Position} + \mathbf{V}(x_0, y_0)\Delta t$.
1. Evaluate velocity: $\mathbf{V}(2, 3) = \langle (3)^2, -(2) \rangle = \langle 9, -2 \rangle$.
2. Calculate displacement: $\Delta \mathbf{r} = \langle 9, -2 \rangle \times 0.1 = \langle 0.9, -0.2 \rangle$.
3. Find new position: $(2, 3) + \langle 0.9, -0.2 \rangle = (2.9, 2.8)$.
\textbf{Final Answer:} $(2.9, 2.8)$.

\vspace{0.3cm}
\textbf{Problem 20.} A particle is located at $(-1, 4)$ moving in $\mathbf{V}(x,y) = \langle 2x+y, xy \rangle$. Estimate its position after a time step of $\Delta t = 0.05$. \\
\textbf{Solution:}
1. Evaluate velocity at $(-1, 4)$: $\mathbf{V}(-1, 4) = \langle 2(-1)+4, (-1)(4) \rangle = \langle 2, -4 \rangle$.
2. Displacement: $\Delta \mathbf{r} = \langle 2, -4 \rangle \times 0.05 = \langle 0.1, -0.2 \rangle$.
3. New position: $(-1, 4) + \langle 0.1, -0.2 \rangle = (-0.9, 3.8)$.
\textbf{Final Answer:} $(-0.9, 3.8)$.

\vspace{0.3cm}
\textbf{Problem 21.} (Quantitative Finance Application) Let risk be defined as $R(w_1, w_2) = 3w_1^2 + 4w_2^2 + 2w_1 w_2$. An algorithm slides down the gradient to minimize risk using step size $\Delta t = 0.01$. If the current weights are $w_1 = 0.4, w_2 = 0.6$, compute the next weights by stepping in the direction of $-\nabla R$. \\
\textbf{Solution:}
1. Find $\nabla R$: $\nabla R = \langle 6w_1 + 2w_2, 8w_2 + 2w_1 \rangle$.
2. Evaluate at $(0.4, 0.6)$: $\nabla R = \langle 6(0.4) + 2(0.6), 8(0.6) + 2(0.4) \rangle = \langle 2.4 + 1.2, 4.8 + 0.8 \rangle = \langle 3.6, 5.6 \rangle$.
3. The displacement is $-\nabla R \times \Delta t = \langle -3.6, -5.6 \rangle \times 0.01 = \langle -0.036, -0.056 \rangle$.
4. New Weights: $(0.4 - 0.036, 0.6 - 0.056) = (0.364, 0.544)$.
\textbf{Final Answer:} The new portfolio weights are $(0.364, 0.544)$.

\vspace{0.3cm}
\textbf{Problem 22.} A drone moves in 3D wind velocity field $\mathbf{V}(x,y,z) = \langle yz, xz, xy \rangle$. If the drone is at $(0, 1, 2)$ at $t=5$, estimate its position at $t=5.2$. \\
\textbf{Solution:}
1. Time step: $\Delta t = 5.2 - 5 = 0.2$.
2. Evaluate velocity: $\mathbf{V}(0, 1, 2) = \langle (1)(2), (0)(2), (0)(1) \rangle = \langle 2, 0, 0 \rangle$.
3. Displacement: $\Delta \mathbf{r} = \langle 2, 0, 0 \rangle \times 0.2 = \langle 0.4, 0, 0 \rangle$.
4. New position: $(0, 1, 2) + \langle 0.4, 0, 0 \rangle = (0.4, 1, 2)$.
\textbf{Final Answer:} $(0.4, 1, 2)$.

\subsection*{Part 5: Dot Product Geometry}

\textbf{Problem 23.} Show mathematically that the vector field $\mathbf{F}(x, y) = \langle y, -x \rangle$ is strictly perpendicular to the radial position vector $\mathbf{r} = \langle x, y \rangle$ everywhere in space. \\
\textbf{Solution:}
Compute the dot product $\mathbf{r} \cdot \mathbf{F}$:
$\langle x, y \rangle \cdot \langle y, -x \rangle = x(y) + y(-x) = xy - xy = 0$.
Since the dot product is exactly 0, the vectors are perfectly orthogonal.
\textbf{Final Answer:} The dot product $\mathbf{r} \cdot \mathbf{F} = 0$, proving it is a purely rotational field.

\vspace{0.3cm}
\textbf{Problem 24.} Does the vector field $\mathbf{F}(x,y) = \langle 5x, 5y \rangle$ model rotational flow or radial flow? Prove using geometry. \\
\textbf{Solution:}
Notice that $\mathbf{F}(x,y) = 5\langle x, y \rangle = 5\mathbf{r}$. 
Because the vector field is just a scalar multiple of the position vector, they point in the exact same direction. 
\textbf{Final Answer:} Radial flow. The vectors are parallel to the position vector (exploding outward).

\subsection*{Part 6: Conservative Vector Fields \& Cross-Partial Test}

\textbf{Problem 25.} Use the cross-partial test to determine if $\mathbf{F}(x,y) = \langle 4x^3 y, x^4 \rangle$ is a conservative vector field. \\
\textbf{Solution:}
Identify $P = 4x^3 y$ and $Q = x^4$.
1. Differentiate $P$ with respect to $y$: $\frac{\partial P}{\partial y} = 4x^3(1) = 4x^3$.
2. Differentiate $Q$ with respect to $x$: $\frac{\partial Q}{\partial x} = 4x^3$.
3. Compare: $4x^3 = 4x^3$. The cross-partials are equal.
\textbf{Final Answer:} Yes, the vector field is conservative.

\vspace{0.3cm}
\textbf{Problem 26.} Determine if the vector field $\mathbf{F}(x,y) = \langle y^2 \sin(x), 2y \cos(x) \rangle$ is conservative. \\
\textbf{Solution:}
Identify $P = y^2 \sin(x)$ and $Q = 2y \cos(x)$.
1. Differentiate $P$ w.r.t $y$: $P_y = 2y \sin(x)$.
2. Differentiate $Q$ w.r.t $x$: $Q_x = 2y(-\sin(x)) = -2y \sin(x)$.
3. Compare: $2y \sin(x) \neq -2y \sin(x)$.
\textbf{Final Answer:} No, the vector field is NOT conservative.

\vspace{0.3cm}
\textbf{Problem 27.} Is the vector field $\mathbf{F}(x,y) = \langle 6x^2 y^2, 4x^3 y \rangle$ conservative? \\
\textbf{Solution:}
Identify $P = 6x^2 y^2$ and $Q = 4x^3 y$.
$P_y = 12x^2 y$.
$Q_x = 12x^2 y$.
Because $P_y = Q_x$, the field is conservative.
\textbf{Final Answer:} Yes.

\vspace{0.3cm}
\textbf{Problem 28.} Evaluate if $\mathbf{F}(x,y) = \langle e^x \cos(y), -e^x \sin(y) \rangle$ is the gradient of a scalar function. \\
\textbf{Solution:}
Being the gradient of a scalar function is synonymous with being a conservative field.
$P_y = \frac{\partial}{\partial y}(e^x \cos(y)) = -e^x \sin(y)$.
$Q_x = \frac{\partial}{\partial x}(-e^x \sin(y)) = -e^x \sin(y)$.
Since $P_y = Q_x$, the field is conservative and must be the gradient of some function $f$.
\textbf{Final Answer:} Yes.

\subsection*{Part 7: "Find the Flaw" - Advanced Error Diagnostics}

\textbf{Problem 29: Partial Derivative Flaw} \\
\textbf{Question:} Find the gradient of $f(x,y) = y^2 \sin(x)$. \\
\textbf{Flawed Solution:} 
1. $f_x = y^2 \cos(x)$. 
2. $f_y = 2y \cos(x)$. 
3. $\nabla f = \langle y^2 \cos(x), 2y \cos(x) \rangle$. \\
\textbf{Your Task:} Find the flaw, explain it, and fix it. \\
\textbf{Correction:} 
The flaw is in Step 2. The student incorrectly applied the derivative of $\sin(x)$ to the $f_y$ calculation. When taking the partial derivative with respect to $y$, the term $\sin(x)$ is treated as a constant.
\textit{Correct Step:} $f_y = 2y \cdot \sin(x)$.
\textit{Correct Final Answer:} $\nabla f = \langle y^2 \cos(x), 2y \sin(x) \rangle$.

\vspace{0.3cm}
\textbf{Problem 30: Cross-Partial Test Flaw} \\
\textbf{Question:} Determine if $\mathbf{F}(x,y) = \langle x^2 y, x y^2 \rangle$ is a conservative field. \\
\textbf{Flawed Solution:}
1. Identify $P = x^2 y$ and $Q = x y^2$.
2. Find $P_x = 2xy$.
3. Find $Q_y = 2xy$.
4. Since $P_x = Q_y$, the vector field is conservative. \\
\textbf{Your Task:} Find the flaw, explain it, and fix it. \\
\textbf{Correction:}
The flaw is a fundamental misunderstanding of the cross-partial test. The test requires taking the derivative of $P$ with respect to \textbf{y} and $Q$ with respect to \textbf{x}, not taking their own respective variables' derivatives.
\textit{Correct Step:} $P_y = x^2$ and $Q_x = y^2$. 
\textit{Correct Final Answer:} Since $x^2 \neq y^2$, the cross-partials are not equal, so the vector field is NOT conservative.

\newpage
\section*{Concept Mapping Index}
Below is a map connecting the core concepts to the problem numbers designed to test them. Use this to review specific weaknesses.

\begin{table}[ht]
\centering
\begin{tabular}{@{}ll@{}}
\toprule
\textbf{Concept / Problem Type} & \textbf{Relevant Question Numbers} \\ \midrule
1. Direct Evaluation of Vector Fields & 1, 2, 3 \\
2. Vector Field Sketching \& Matching (2D) & 4, 5, 6, 7, 8 \\
3. Vector Field Sketching \& Matching (3D) & 9, 10, 11 \\
4. Gradient Vector Fields \& Partial Derivatives & 12, 13, 14, 15, 16, 17, 18 \\
5. Linear Approximation (Euler's Method) \& Kinematics & 19, 20, 21, 22 \\
6. Dot Product for Geometry \& Orthogonality & 23, 24 \\
7. Conservative Fields \& Cross-Partial Test & 25, 26, 27, 28 \\
8. Diagnostic Error Analysis ("Find the Flaw") & 29, 30 \\ \bottomrule
\end{tabular}
\end{table}

\end{document}